% Fragmento público generado desde propietarios íntegros.
% Residencia: 52.6.1.
% Fuente: ../incoming/radion_monografia_20260827/manuscrito/sections/07_incertidumbre_oscilador.tex:1-111
\noindent La forma simpléctica coordina acción, incertidumbre y dinámica oscilatoria sin alterar la genealogía del estado.\par\medskip

\subsubsection{Realización del estado en un espacio de Hilbert}

\paragraph{Matriz de covarianza mínima}

Una vez construida la elipse de acción, puede realizarse en un par de
operadores conjugados
\[
[\widehat Q,\widehat P]=i\hret.
\]
Definimos
\begin{equation}
V_\eta=\frac{\hret}{2}
\begin{pmatrix}
e^\eta&0\\[2pt]0&e^{-\eta}
\end{pmatrix}.
\label{eq:covariance}
\end{equation}
Su determinante es
\begin{equation}
\det V_\eta=\frac{\hret^2}{4}.
\label{eq:cov-det}
\end{equation}
Por consiguiente,
\begin{equation}
(\Delta Q)^2=\frac{\hret}{2}e^\eta,
\qquad
(\Delta P)^2=\frac{\hret}{2}e^{-\eta},
\qquad
\Delta Q\,\Delta P=\frac{\hret}{2}.
\label{eq:uncertainty}
\end{equation}
La desigualdad de Robertson--Schrödinger se satura.

El aniquilador comprimido
\begin{equation}
\widehat a_\eta=
\frac{e^{-\eta/2}\widehat Q+i e^{\eta/2}\widehat P}
{\sqrt{2\hret}}
\label{eq:a-eta}
\end{equation}
satisface \([\widehat a_\eta,\widehat a_\eta^\dagger]=1\). Su vacío tiene
función de onda
\begin{equation}
\psi_\eta(Q)=
(\pi\hret e^\eta)^{-1/4}
\exp\left(-\frac{Q^2}{2\hret e^\eta}\right).
\label{eq:gaussian}
\end{equation}

\paragraph{Elipse como nivel del funcional de Mahalanobis}

El conjunto
\begin{equation}
\mathbb E_{\mathrm{act}}
=\{x\in\R^2:x^{\mathsf T}V_\eta^{-1}x\le4\}
\label{eq:mahalanobis}
\end{equation}
tiene semiejes
\begin{equation}
a_S=2\Delta Q,\qquad b_S=2\Delta P.
\label{eq:two-sigma}
\end{equation}
Su área es
\[
4\pi\sqrt{\det V_\eta}=2\pi\hret=\hfull.
\]
La elipse HMT construida antes coincide exactamente con el contorno de dos
desviaciones del estado gaussiano mínimo.

\begin{proposition}[Cadena única del cociente de forma]
Tras declarar las interfaces espectral, simpléctica y de Hilbert,
\begin{equation}
\boxed{
e^{-\eta}
=\sqrt{\frac{A-C^*}{A+C^*}}
=\frac{b_1}{a_1}
=\frac{b_S}{a_S}
=\frac{\Delta P}{\Delta Q}.}
\label{eq:shape-chain-hilbert}
\end{equation}
\end{proposition}

La elipse de acción corresponde al nivel clásico \(H=\hret\omega\) y al
contorno \(2\sigma\) del vacío comprimido. No debe confundirse con la energía
media del vacío, \(\hret\omega/2\). Esta diferencia de factor dos es un
control físico obligatorio.

\paragraph{Condiciones de invariancia de la elipse como objeto operatorio}

En la formulación convencional, una elipse de incertidumbre puede verse como
un recurso gráfico para representar una matriz de covarianza. Aquí el orden
causal es inverso: la elipse positiva ya existe como publicación de
\((A,C^*,\hret)\); el espacio de Hilbert la reconoce mediante
\eqref{eq:covariance}. La realidad ontológica reclamada no significa que una
partícula clásica recorra físicamente un óvalo en el plano \((Q,P)\). Significa
que el estado conserva una forma positiva, una acción total y un cociente de
orientación antes de que una teoría probabilista los interprete.

La función de onda no crea la forma; la representa. La incertidumbre no crea
\(\hbar\); publica la imposibilidad de comprimir simultáneamente las dos
direcciones sin alterar el producto simpléctico.

\begin{narrativebox}
La lectura probabilista es una realización posible de la elipse, no su causa.
El objeto primario es la conservación conjunta de producto y orientación. La
probabilidad aparece cuando esa geometría se publica como covarianza de un
estado sobre Hilbert.
\end{narrativebox}

